%================================================================
\chapter{Python Environment and Reproducibility}
\label{app:env}
%================================================================

The following packages and versions were used in this study. A complete \texttt{requirements.txt} file is available in the project GitHub repository.

\begin{lstlisting}[language={},caption={Core Python dependencies (requirements.txt).},label={lst:requirements}]
pandas==2.0.3
numpy==1.24.4
nltk==3.8.1
spacy==3.6.1
scikit-learn==1.3.0
torch==2.0.1
transformers==4.30.2
vaderSentiment==3.3.2
imbalanced-learn==0.11.0
shap==0.42.1
lime==0.2.0.1
emoji==2.8.0
matplotlib==3.7.2
seaborn==0.12.2
scipy==1.11.1
joblib==1.3.1
\end{lstlisting}

%================================================================
\chapter{Dataset Exploratory Data Analysis Outputs}
\label{app:eda}
%================================================================

This appendix contains supplementary \gls{eda} outputs. Table~\ref{tab:eda_summary} summarises key corpus statistics.

\begin{table}[H]
\centering
\caption{Sentiment140 corpus summary statistics.}
\label{tab:eda_summary}
\begin{tabular}{lc}
  \toprule
  \rowcolor{uwsblue!90}
  \textcolor{white}{\textbf{Statistic}} & \textcolor{white}{\textbf{Value}} \\
  \midrule
  \rowcolor{tablegrey}
  Total raw records            & 1,600,000 \\
  Duplicate records removed    & 429 \\
  \rowcolor{tablegrey}
  Usable records               & 1,599,571 \\
  Positive class (\%)          & 49.97\% \\
  \rowcolor{tablegrey}
  Negative class (\%)          & 49.92\% \\
  Neutral class (\%)           & 0.11\% \\
  \rowcolor{tablegrey}
  Median tweet length (chars)  & 74 \\
  Mean tweet length (chars)    & 73.6 \\
  \rowcolor{tablegrey}
  Temporal coverage            & April -- June 2009 \\
  Unique users                 & 659,775 \\
  \rowcolor{tablegrey}
  TF-IDF vocabulary (cap)      & 10,000 \\
  \bottomrule
\end{tabular}
\end{table}

%================================================================
\chapter{ROC-AUC Computation Code}
\label{app:roc}
%================================================================

Listing~\ref{lst:rocauc} presents the complete \gls{roc}-\gls{auc} evaluation code used in this study.

\begin{lstlisting}[language=Python, caption={Multi-class ROC-AUC evaluation (One-vs-Rest strategy).}, label={lst:rocauc}]
from sklearn.metrics import (roc_auc_score, roc_curve,
                              auc, classification_report)
from sklearn.preprocessing import label_binarize
import matplotlib.pyplot as plt
import numpy as np

# Map labels to 0,1,2 for sklearn
label_map = {-1: 0, 3: 1, 5: 2}
y_test_m   = y_test.map(label_map).values
classes    = [0, 1, 2]
class_names = ['Negative', 'Neutral', 'Positive']
y_test_bin  = label_binarize(y_test_m, classes=classes)

results = {}
for name, model in trained_models.items():
    y_proba = model.predict_proba(X_test)

    # Per-class AUC (One-vs-Rest)
    auc_per_class = roc_auc_score(y_test_bin, y_proba,
                                  multi_class='ovr',
                                  average=None)
    auc_macro   = roc_auc_score(y_test_bin, y_proba,
                                multi_class='ovr',
                                average='macro')
    auc_weighted = roc_auc_score(y_test_bin, y_proba,
                                 multi_class='ovr',
                                 average='weighted')
    results[name] = {
        'AUC_Neg':      auc_per_class[0],
        'AUC_Neu':      auc_per_class[1],
        'AUC_Pos':      auc_per_class[2],
        'Macro AUC':    auc_macro,
        'Weighted AUC': auc_weighted,
    }

# Plot ROC curves
colors = ['#1F3864', '#2E75B6', '#70AD47']
fig, axes = plt.subplots(1, len(trained_models),
                         figsize=(20, 5))
for ax, (name, model) in zip(axes, trained_models.items()):
    y_proba = model.predict_proba(X_test)
    for i, (cls, col) in enumerate(
            zip(class_names, colors)):
        fpr, tpr, _ = roc_curve(y_test_bin[:, i],
                                 y_proba[:, i])
        ax.plot(fpr, tpr, color=col, lw=2,
                label=f'{cls} (AUC={auc(fpr,tpr):.3f})')
    ax.plot([0,1],[0,1],'k--',lw=1,label='Random')
    ax.set_title(f'ROC — {name}', fontsize=11)
    ax.set_xlabel('False Positive Rate')
    ax.set_ylabel('True Positive Rate')
    ax.legend(loc='lower right', fontsize=8)
    ax.grid(alpha=0.3)
plt.tight_layout()
plt.savefig('results/roc_curves_all_models.pdf', dpi=200)
\end{lstlisting}

%================================================================
\chapter{Source Code Notebooks}
\label{app:notebooks}
%================================================================

The complete experimental pipeline is implemented across six Google Colab notebooks. Each notebook is self-contained and designed to be executed sequentially. Table~\ref{tab:notebooks} summarises the notebooks and their scope.

\begin{table}[H]
\centering
\caption{Google Colab source code notebooks for the experimental pipeline.}
\label{tab:notebooks}
\begin{tabular}{@{}clp{7cm}@{}}
  \toprule
  \rowcolor{uwsblue!90}
  \textcolor{white}{\textbf{No.}} &
  \textcolor{white}{\textbf{Notebook}} &
  \textcolor{white}{\textbf{Description}} \\
  \midrule
  \rowcolor{tablegrey}
  01 & \texttt{01\_EDA\_and\_Preprocessing.ipynb} & Loading the Sentiment140 dataset, exploratory data analysis (class distribution, tweet length, temporal volume), full 7-step text preprocessing pipeline, structural feature extraction, and export to Parquet format. \\
  02 & \texttt{02\_Feature\_Engineering.ipynb} & TF-IDF vectorisation (10,000-feature cap, unigram+bigram), VADER compound score extraction, structural feature assembly, and SMOTE oversampling for class balance. \\
  \rowcolor{tablegrey}
  03 & \texttt{03\_Model\_Training.ipynb} & Training of Logistic Regression, Random Forest, MLP Neural Network, and Voting Ensemble models; hyperparameter configuration; early stopping for MLP; model serialisation with \texttt{joblib}. \\
  04 & \texttt{04\_Temporal\_Analysis.ipynb} & Daily sentiment aggregation, 7-day rolling mean smoothing, sentiment velocity computation, Pearson lag correlation analysis ($k = 1$--$7$ days), and temporal dynamics visualisation. \\
  \rowcolor{tablegrey}
  05 & \texttt{05\_Evaluation\_ROC\_AUC.ipynb} & Held-out test set evaluation; classification report (accuracy, precision, recall, F1); multi-class ROC-AUC computation (One-vs-Rest); normalised confusion matrix generation. \\
  06 & \texttt{06\_XAI\_SHAP\_LIME.ipynb} & SHAP TreeExplainer global and local feature importance analysis for Random Forest; LIME text explanation generation for Logistic Regression; XAI audit and visualisation. \\
  \bottomrule
\end{tabular}
\end{table}

All notebooks were developed and executed in Google Colab (Python 3.10) using the dependency versions listed in Appendix~\ref{app:env}. The notebooks follow good scientific computing practices, including explicit random seed setting (\texttt{RANDOM\_STATE = 42}), modular function definitions, and inline commentary \citep{wilson2017good}. Results are exported to a shared \texttt{/results/} directory in Google Drive for reproducibility.
